\documentclass[12pt]{article}
\usepackage[utf8]{inputenc}
\usepackage[T1]{fontenc}
\usepackage{hyperref}
\title{West Damar Language Resources for K-12 Curriculum Development}
\author{Compiled by Claude Sonnet 5.5 (claude-sonnet-5-5) for Andr\'as Kornai}
\date{}
\begin{document}
\maketitle

\section{Introduction}
West Damar (ISO 639-3 \texttt{drn}; Glottocode west2548), also called North Damar or Damar Batumerah, is an Austronesian language spoken in two villages on the northwest side of Damar Island, Maluku, Indonesia. Wikipedia cites about 800 speakers (1987). Its classification is debated. Very little material exists beyond linguistic description.

\section{Primary Educational Resources}
No primers, textbooks or school programmes were found in the search results.

\section{Linguistic and Reference Materials}
\begin{itemize}
\item Svetlana F.\ Chlenova (2008), preliminary grammatical notes on Damar Batumerah (West Damar), Glottolog entry: \url{https://glottolog.org/resource/reference/id/159808}
\item Chlenova, SEALS XVI abstract: \url{https://lingweb.eva.mpg.de/archive/jakarta/seals/Chlenova_SEALS_XVI_Abstract.pdf}
\item Chlenov and Chlenova (2008), Glottolog entry: \url{https://glottolog.org/resource/reference/id/42376}
\item Grambank source entry for Chlenova's West Damar description: \url{https://grambank.clld.org/sources/sChlenovaWDamar}
\item Wikipedia overview: \url{https://en.wikipedia.org/wiki/West_Damar_language}
\end{itemize}

\section{Religious and Community Resources}
None found.

\section{Digital and Audio Resources}
None found.

\section{Curriculum Development Context}
The language is not known to be taught in schools. The documentation found is linguistic and preliminary, with no orthography, wordlist or teaching material in the results. Any curriculum would require new documentation with the community. Some speakers were reportedly relocated to Seram in the late 1960s to 1980s, which may affect the speaker community.

\end{document}